% 実行モデル詳細
\section{実行モデル詳細}

\subsection{ライフサイクル}

本アーキテクチャのモジュールは、以下のライフサイクルに従う：

\begin{enumerate}
    \item \textbf{インスタンス生成}（グローバルスコープ） \\
    モジュールのコンストラクタが呼ばれる。ModuleTimerのメンバ変数が\texttt{\_startTime = 0}で初期化される。
    この時点では\texttt{millis()}は信頼できないため、タイマーの実質的な初期化は行わない。

    \item \textbf{init()呼び出し}（\texttt{setup()}内） \\
    ハードウェアの初期化、ピンモード設定、ModuleTimerの\texttt{setTime()}呼び出しを行う。
    失敗時は\texttt{false}を返し、呼び出し側でエラーハンドリングを行う。

    \item \textbf{update()ループ}（\texttt{loop()}内） \\
    毎ループ呼び出される。ModuleTimerの\texttt{getNowTime()}で経過時間を確認し、
    必要なタイミングでのみ実処理を実行する。
\end{enumerate}

\subsection{タイミング制御}

\texttt{delay()}は一切使用しない。すべてのタイミング制御はModuleTimerで行う。

\begin{tabularx}{\textwidth}{l l X}
\toprule
\textbf{パターン} & \textbf{方法} & \textbf{コード例} \\
\midrule
周期実行 &
\texttt{getNowTime() >= interval}で判定し、\texttt{setTime()}でリセット &
1秒ごとのセンサー読み取り \\
デバウンス &
入力変化時に\texttt{setTime()}でリセット、安定後に確定 &
ボタンのチャタリング除去 \\
ワンショット &
\texttt{getNowTime()}で経過時間を確認、フラグで制御 &
起動後N秒でLED消灯 \\
タイムアウト &
\texttt{getNowTime() >= timeout}で時間超過を検出 &
WiFi接続のタイムアウト \\
\bottomrule
\end{tabularx}

\subsection{モジュール間データフロー}

\texttt{SystemData}を介したデータフローの例：

\begin{lstlisting}[style=cpp, caption={モジュール間データフローの例}]
// 温度センサーモジュール（先に実行）
void TempSensorModule::update(SystemData* sys) {
    if (readTimer.getNowTime() >= 1000) {
        readTimer.setTime();
        sys->temp.temperature = readSensor();
        sys->temp.isValid = true;
    }
}

// LEDモジュール（後に実行）
void LedModule::update(SystemData* sys) {
    // 温度データを読み取って、温度に応じてLEDを制御
    if (sys->temp.isValid && sys->temp.temperature > 30.0) {
        sys->led.state = true;   // 30度超でLED点灯
    } else {
        sys->led.state = false;
    }
    digitalWrite(LED_PIN, sys->led.state);
}
\end{lstlisting}

\subsection{モジュール実装テンプレート}

新規モジュールを追加する際の標準テンプレート：

\begin{lstlisting}[style=cpp, caption={モジュール実装テンプレート}]
// ExampleModule.h
#pragma once
#include "IModule.h"
#include "ModuleTimer.h"

// このモジュール固有のデータ構造体
struct ExampleData {
    // モジュールが管理するデータを定義
    int value;
    bool hasError;
};

class ExampleModule : public IModule {
private:
    ModuleTimer _updateTimer;  // メンバ変数としてタイマー保持

public:
    ExampleModule() {}  // コンストラクタ

    bool init() override {
        _updateTimer.setTime();  // タイマー初期化
        // ハードウェア初期化処理...
        return true;
    }

    void update(SystemData* sys) override {
        if (_updateTimer.getNowTime() >= 100) {
            _updateTimer.setTime();
            // 更新処理...
        }
    }
};
\end{lstlisting}
